\documentclass[11pt,a4paper]{article}
\usepackage{gtu-el-yazisi}
\begin{document}
\notbilgi{Nümerik Analiz}{11. Hafta -- Lineer denklem sistemleri için iterasyon yöntemleri}{numerik-anlaiz-11.hafta-ders-notu-ve-uygulamasi.pdf}

\sayfa{1}
\begin{center}
  {\Large\bfseries MAT 214}\\[2mm]
  Sayısal Yöntemler\\[4mm]
  \underline{11. Hafta}\\[6mm]
  \textbf{Lineer Cebri Denklemler Sistemi için\\ İterasyon Yöntemleri.}
\end{center}

\begin{enumerate}[label=\arabic*)]
  \item Matrisin öz değer ve öz fonksiyonları;
  \item Jacobi iterasyon yöntemi.
\end{enumerate}

\sayfa{2}
\section{Lineer Cebri Denklem Sistemleri için İterasyon Yöntemleri}
\subsection{Matrisin özdeğer ve özvektörleri, Spektral Yarıçap.}
$n\times n$ boyutlu bir $A=[a_{ij}]$ matrisi verilmiş olsun. Bu matris için, $\lambda$ reel veya kompleks sayı olmakla
\[ \underline{Ax=\lambda x} \ \Leftrightarrow\ (A-\lambda I_n)\,x=0, \qquad I_n \text{ özdeşlik matristir} \]
denklem sistemini inceleyelim. Bu denklem sisteminin açık şekli
\[
  \begin{cases}
    \underline{(a_{11}-\lambda)}\,x_1+a_{12}x_2+\cdots+a_{1n}x_n=0,\\
    a_{21}x_1+\underline{(a_{22}-\lambda)}\,x_2+\cdots+a_{2n}x_n=0,\\
    \qquad\vdots\qquad\qquad\qquad\qquad\vdots\\
    a_{n1}x_1+a_{n2}x_2+\cdots+\underline{(a_{nn}-\lambda)}\,x_n=0
  \end{cases}
  \qquad x=\begin{bmatrix}x_1\\x_2\\\vdots\\x_n\end{bmatrix}
\]
olur. Kolayca görüldüğü üzere her $\lambda$ sayısı için $x=0$ bu sistemin çözümü olur. Bu çözüme
\underline{aşikar} (\underline{trivial}) çözüm denir. Maksat bu denklem sisteminin trivial olmayan
($x\neq 0$) çözümünü bulmaktır.

$\lambda$ parametresinin, bu sistemin \underline{trivial olmayan} çözümlerini temin eden değerlerine
\underline{özdeğerler}, trivial olmayan çözümlere ise \underline{özvektörler} denir.

Kolayca görüldüğü üzere, eğer $\det(A-\lambda I_n)\neq 0$ olursa sistemin tek $x=0$ trivial çözümü olur.
Demek ki trivial olmayan çözüme $\lambda$'nın $\det(A-\lambda I_n)=0$ yapan değerleri ile ulaşabiliriz.
Başka deyimle,
\[ p(\lambda)=\det(A-\lambda I)\text{'ye } A \text{ matrisinin karakteristik} \]

\sayfa{3}
polinomu dersek, bu polinomun sıfırlarına $A$ matrisinin özdeğerleri denir:
\[ p(\lambda)=0 \ \Leftrightarrow\ \underline{\det(A-\lambda I)=0} \]
denklemlerinin çözümleri $A$ matrisinin özdeğerleridir.
Eğer $\lambda$ sayısı $A$ matrisinin bir özdeğeri ise
\[ (A-\lambda I)\,x=0 \]
sisteminin trivial olmayan ($x\neq 0$) çözümleri bu özdeğere karşı gelen özvektörler olur.

Böylece $n\times n$ boyutlu $A=[a_{ij}]$ matrisinin özdeğer ve özfonksiyonlarını aşağıdaki aşamalarla bulabiliriz:
\begin{itemize}
  \item $\det(A-\lambda I)=0$ denklemini çözerek, $A$ matrisinin $\lambda$ \underline{özdeğerleri} bulunur;
  \item $\lambda$ özdeğer olduğunda
  \[ (A-\lambda I)\,x=0 \]
  denklem sistemini çözmekle, $\lambda$ özdeğerine karşı gelen $x\neq 0$ \underline{özvektörleri} bulunur;
\end{itemize}

\begin{ornek}
Aşağıdaki matrislerin özdeğer ve özvektörlerini bulunuz\\[1mm]
a) \ $A=\begin{bmatrix}0&1\\-1&0\end{bmatrix}$, \qquad b) \ $A=\begin{bmatrix}2&0&0\\1&1&2\\1&-1&4\end{bmatrix}$
\end{ornek}
\begin{cozum}
\[ \text{a)}\quad A-\lambda I=\begin{bmatrix}0&1\\-1&0\end{bmatrix}-\lambda\begin{bmatrix}1&0\\0&1\end{bmatrix}=\begin{bmatrix}-\lambda&1\\-1&-\lambda\end{bmatrix} \]

\sayfa{4}
\[ \det(A-\lambda I)=\begin{vmatrix}-\lambda&1\\-1&-\lambda\end{vmatrix}=\lambda^2+1=0 \]
\[ \Rightarrow\ \lambda_1=i,\ \lambda_2=-i \]
$A$ matrisinin iki $\lambda_1=i$ ve $\lambda_2=-i$ özdeğerleri vardır.
\[
  \lambda_1=i \ \rightarrow\ (A-\lambda_1 I)\,x=\begin{bmatrix}-i&1\\-1&-i\end{bmatrix}\begin{bmatrix}x_1\\x_2\end{bmatrix}=0
\]
\[
  \Rightarrow\ \begin{cases}-i\,x_1+x_2=0\\-x_1-i\,x_2=0\end{cases}
  \ \Rightarrow\ \underline{x_2=i\,x_1}
\]
\[
  \begin{bmatrix}x_1\\x_2\end{bmatrix}=\begin{bmatrix}x_1\\i\,x_1\end{bmatrix}
  =\underline{\begin{bmatrix}1\\i\end{bmatrix}}\,x_1
\]
Böylece, \underline{$\lambda_1=i$} özdeğerine karşı gelen özvektör \underline{$\begin{bmatrix}1\\i\end{bmatrix}$}'dir
\[
  \lambda_2=-i \ \rightarrow\ (A-\lambda_2 I)\,x=\begin{bmatrix}i&1\\-1&i\end{bmatrix}\begin{bmatrix}x_1\\x_2\end{bmatrix}=0
\]
\[
  \Rightarrow\ \begin{cases}i\,x_1+x_2=0\\-x_1+i\,x_2=0\end{cases}
  \ \rightarrow\ x_2=-i\,x_1
\]
\[
  \begin{bmatrix}x_1\\x_2\end{bmatrix}=\begin{bmatrix}x_1\\-i\,x_1\end{bmatrix}
  =\underline{\begin{bmatrix}1\\-i\end{bmatrix}}\,x_1
\]
\underline{$\lambda_2=-i$} özdeğerine karşı gelen özvektör \underline{$\begin{bmatrix}1\\-i\end{bmatrix}$}'dir
\[
  \text{b)}\quad A-\lambda I=\begin{bmatrix}2&0&0\\1&1&2\\1&-1&4\end{bmatrix}
  -\lambda\begin{bmatrix}1&0&0\\0&1&0\\0&0&1\end{bmatrix}
  =\begin{bmatrix}2-\lambda&0&0\\1&1-\lambda&2\\1&-1&4-\lambda\end{bmatrix}
\]
\begin{align*}
  \det(A-\lambda I)&=\begin{vmatrix}2-\lambda&0&0\\1&1-\lambda&2\\1&-1&4-\lambda\end{vmatrix}
  =(2-\lambda)\cdot\begin{vmatrix}1-\lambda&2\\-1&4-\lambda\end{vmatrix}\\
  &=(2-\lambda)(6-5\lambda+\lambda^2)=0
\end{align*}

\sayfa{5}
\[ \Rightarrow\ \lambda_{1,2}=2,\ \lambda_3=3. \]
$A$ matrisinin birisi iki katlı olmakla 2 sayıda $\lambda_{1,2}=2,\ \lambda_3=3$ özdeğeri vardır.
\[
  \lambda_{1,2}=2 \ \rightarrow\ (A-\lambda_{1,2}I)\,x=0 \ \Rightarrow\
  \begin{bmatrix}0&0&0\\1&-1&2\\1&-1&2\end{bmatrix}\begin{bmatrix}x_1\\x_2\\x_3\end{bmatrix}=0
\]
\[
  \Rightarrow\ x_1-x_2+2x_3=0 \ \Rightarrow\ x_1=x_2-3x_3\stepcounter{mpfootnote}\textsuperscript{\thempfootnote}
\]
\footnotetext{Claude notu: $x_1-x_2+2x_3=0$ denkleminden muhtemelen $x_1=x_2-2x_3$ olmalı (bu durumda ikinci özvektör $[-2,\,0,\,1]^T$ olur).}
\[
  \begin{bmatrix}x_1\\x_2\\x_3\end{bmatrix}=\begin{bmatrix}x_2-3x_3\\x_2\\x_3\end{bmatrix}
  =\underline{\begin{bmatrix}1\\1\\0\end{bmatrix}}\,x_2+\underline{\begin{bmatrix}-3\\0\\1\end{bmatrix}}\,x_3
\]
\underline{$\lambda_{1,2}$}$=2$ özdeğerine karşı gelen lineer bağımsız iki tane özvektör
\[
  \underline{\begin{bmatrix}1\\1\\0\end{bmatrix}}\ \text{ve}\ \underline{\begin{bmatrix}-3\\0\\1\end{bmatrix}}\,\text{'dur.}
\]
\[
  \lambda_3=3 \ \rightarrow\ (A-\lambda_3 I)\,x=0 \ \Rightarrow\
  \begin{bmatrix}-1&0&0\\0&-2&2\\1&-1&1\end{bmatrix}\begin{bmatrix}x_1\\x_2\\x_3\end{bmatrix}=0\stepcounter{mpfootnote}\textsuperscript{\thempfootnote}
\]
\footnotetext{Claude notu: $A-3I$ matrisinin ikinci satırı muhtemelen $[1,\,-2,\,2]$ olmalı (ikinci denklem $x_1-2x_2+2x_3=0$ olurdu; $x_1=0$ olduğu için sonuç değişmez).}
\[
  \Rightarrow\ \begin{cases}-x_1=0\\-2x_2+2x_3=0\\x_1-x_2+x_3=0\end{cases}
  \ \Rightarrow\ x_1=0,\ x_2=x_3
\]
\[
  \begin{bmatrix}x_1\\x_2\\x_3\end{bmatrix}=\begin{bmatrix}0\\x_2\\x_2\end{bmatrix}
  =\underline{\begin{bmatrix}0\\1\\1\end{bmatrix}}\,x_2
\]
$\lambda_3=3$ özdeğerine karşı gelen özvektör $\begin{bmatrix}0\\1\\1\end{bmatrix}$'dir.

\end{cozum}

\sayfa{6}
$n\times n$ boyutlu $A$ matrisinin
\[ \lambda_1,\lambda_2,\dots,\lambda_m \]
\noindent\begin{minipage}[b]{0.55\linewidth}
özdeğerleri bulunsun.
\end{minipage}\hfill
\begin{minipage}[b]{0.38\linewidth}
$\lambda=a+ib$\\[1mm]
$|\lambda|=\sqrt{a^2+b^2}$
\end{minipage}
\[ r(A)=\max\bigl\{|\lambda_1|,\dots,|\lambda_m|\bigr\}, \]
sayısına $A$ matrisinin \underline{spektral yarıçap} (radius) denir.

\medskip
\noindent\underline{\underline{Örneğin}} \quad \underline{a)} \ $A=\begin{bmatrix}0&1\\-1&0\end{bmatrix}$ matrisinin özdeğerleri
$\lambda_1=i$ ve $\lambda_2=-i$ olduğuna göre
$|\lambda_1|=1$ ve $|\lambda_2|=1$ olduğundan, spektral yarıçap
\[ r(A)=\max\{1,1\}=1 \quad\text{olur.} \]
\underline{\underline{b)}} \ $A=\begin{bmatrix}2&0&0\\1&1&2\\1&-1&4\end{bmatrix}$ matrisinin özdeğerleri $\lambda_{1,2}=2$ ve $\lambda_3=3$
olduğuna göre $|\lambda_{1,2}|=2$ ve $|\lambda_3|=3$ olduğuna göre Spektral yarıçap
\[ r(A)=\max\{2,3\}=3 \quad\text{olur.} \]

\subsection{Yakınsak Matrisler.}
$n\times n$ boyutlu $A=[a_{ij}]$ matrisi için
\[
  A^2=A\cdot A,\quad A^3=A\cdot A\cdot A,\ \dots,\quad
  A^k=\underbrace{A\cdot A\cdots A}_{k\ \text{kez}}
  \qquad \text{matrisleri düzenleyelim}
\]
$[A^k]_{ij}$ ile $A^k$ matrisinin $i$. satır ve $j$. sütun elemanını işaretleyelim.

\sayfa{7}
Eğer her $i,j=1,2,..,n$ için
\[ \lim_{k\to+\infty}[A^k]_{ij}=0 \]
olursa, diyorlar ki, $A$ matrisi \underline{yakınsak matristir}

\begin{ornek}[2]
$A=\begin{bmatrix}\frac12&0\\[1mm]\frac14&\frac12\end{bmatrix}$ matrisinin yakınsak matris olduğunu gösteriniz.
\end{ornek}
\begin{cozum}
\[
  A^2=A\cdot A=\begin{bmatrix}\frac12&0\\[1mm]\frac14&\frac12\end{bmatrix}\cdot\begin{bmatrix}\frac12&0\\[1mm]\frac14&\frac12\end{bmatrix}
  =\begin{bmatrix}\frac14&0\\[1mm]\frac28&\frac14\end{bmatrix}
  =\begin{bmatrix}\frac14&0\\[1mm]\frac14&\frac14\end{bmatrix}
\]
\[
  A^3=A\cdot A^2=\begin{bmatrix}\frac12&0\\[1mm]\frac14&\frac12\end{bmatrix}\cdot\begin{bmatrix}\frac14&0\\[1mm]\frac14&\frac14\end{bmatrix}
  =\begin{bmatrix}\frac18&0\\[1mm]\frac{3}{16}&\frac18\end{bmatrix}
\]
\[
  A^4=A\cdot A^3=\begin{bmatrix}\frac{1}{16}&0\\[1mm]\frac18&\frac{1}{16}\end{bmatrix}\quad\text{vs.}
\]
\[
  A^k=\begin{bmatrix}\frac{1}{2^k}&0\\[1mm]\frac{k}{2^{k+1}}&\frac{1}{2^k}\end{bmatrix}
\]
Görüldüğü üzere
\[
  \lim_{k\to+\infty}\frac{1}{2^k}=0 \quad\text{ve}\quad \lim_{k\to+\infty}\frac{k}{2^{k+1}}=0
\]
\[ \Rightarrow\ A \text{ yakınsak matristir.} \]
\end{cozum}
$A$ matrisinin \underline{yakınsak olması}, onun spektral yarıçapının
\underline{$r(A)<1$} olması ile \underline{gerek ve yeterdir.}

\sayfa{8}
\section{Lineer Cebri Denklem sistemi için Jacobi iterasyon yöntemi.}
$n\times n$ boyutlu $A=[a_{ij}]$ matrisi verilsin, öyle ki, $a_{ii}\neq 0$ olsun.
$x=[x_1,\dots,x_n]^T$ bilinmeyeni ve $b=[b_1,\dots,b_n]^T$ sağ tarafı için
\[
  Ax=b \ \Leftrightarrow\
  \begin{cases}
    \kutula{a_{11}x_1}+a_{12}x_2+\cdots+a_{1n}x_n=b_1\\
    a_{21}x_1+\kutula{a_{22}x_2}+\cdots+a_{2n}x_n=b_2\\
    \qquad\vdots\qquad\qquad\qquad\qquad\qquad\ \vdots\\
    a_{n1}x_1+a_{n2}x_2+\cdots+\kutula{a_{nn}x_n}=b_n
  \end{cases}
\]
denklem sisteminin çözümü için \okunamadi{yaklaşık} (iteratif) yöntem verelim.

1. denklemin her tarafına $a_{11}$, 2. denklemin her tarafına $a_{22}$ ve s. $n$. denklemin her tarafına $a_{nn}$'le bölürsek
\[
  \begin{array}{r@{\;}c@{\;}l@{\;}c@{\;}l@{\;}c@{\;}l}
    x_1&=&-\dfrac{a_{12}}{a_{11}}\,x_2&-\cdots-&\dfrac{a_{1n}}{a_{11}}\,x_n&+&\dfrac{b_1}{a_{11}}\,,\\[4mm]
    x_2&=&-\dfrac{a_{21}}{a_{22}}\,x_1&-\cdots-&\dfrac{a_{2n}}{a_{22}}\,x_n&+&\dfrac{b_2}{a_{22}}\,,\\[2mm]
    \vdots&&\vdots&&&&\vdots\\[2mm]
    x_n&=&-\dfrac{a_{n1}}{a_{nn}}\,x_1&-\cdots-&\dfrac{a_{n,n-1}}{a_{nn}}\,x_{n-1}&+&\dfrac{b_n}{a_{nn}}
  \end{array}
\]
Şeklinde veya kısaca
\[
  x_i=\sum_{\substack{j=1\\ j\neq i}}^{n}\left(-\frac{a_{ij}}{a_{ii}}\,x_j\right)+\frac{b_i}{a_{ii}}\,,
  \qquad i=1,2,\dots,n
\]
denklem sistemini elde ederiz.

\sayfa{9}
Başlangıç iterasyon olarak
\[ x^{(0)}=0 \qquad \left(x^{(0)}=\begin{bmatrix}0\\0\\\vdots\\0\end{bmatrix}\right) \]
alırsak (veya $x^{(0)}$ olarak keyfi bir vektör de alınabilir)
\[
  x_i^{(k)}=\frac{1}{a_{ii}}\left[-\sum_{\substack{j=1\\ j\neq i}}^{n}a_{ij}\,x_j^{(k-1)}+b_i\right],
  \qquad i=1,2,\dots,n;\quad k=1,2,\dots
\]
Şeklinde iterasyon belirleyebiliriz. Bu iterasyona Jacobi iterasyonu denir. $k$ iterasyon adımıdır.

\begin{ornek}[3]
\[
  \begin{cases}
    \kutula{3x_1}-x_2+x_3=1\\
    3x_1+\kutula{6x_2}+2x_3=0\\
    3x_1+3x_2+\kutula{7x_3}=0
  \end{cases}
\]
denklem sistemi için $x^{0}=0$ olmakla Jacobi iterasyon yöntemi uygulayarak ilk 2 iterasyonu bulunuz
\end{ornek}
\begin{cozum}
1. denklemi 3 ile, 2. denklemi 6 ile ve 3. denklemi 7 ile bölerek sistemi aşağıdaki şekilde yazabiliriz
\[
  \left.\begin{aligned}
    \kutula{x_1}-\tfrac13x_2+\tfrac13x_3&=\tfrac13\\[1mm]
    \tfrac36x_1+\kutula{x_2}+\tfrac26x_3&=0\\[1mm]
    \tfrac37x_1+\tfrac37x_2+\kutula{x_3}&=0
  \end{aligned}\right\}
  \ \Rightarrow\
  \left\{\begin{aligned}
    x_1&=\tfrac13x_2-\tfrac13x_3+\tfrac13\\[1mm]
    x_2&=-\tfrac12x_1-\tfrac13x_3\\[1mm]
    x_3&=-\tfrac37x_1-\tfrac37x_2
  \end{aligned}\right.
\]
Şimdi Jacobi iterasyonunu belirleyelim:

\sayfa{10}
\[
  \left\{\begin{aligned}
    x_1^{(k)}&=\frac13\,x_2^{(k-1)}-\frac13\,x_3^{(k-1)}+\frac13\,,\\[1mm]
    x_2^{(k)}&=-\frac12\,x_1^{(k-1)}-\frac13\,x_3^{(k-1)}\,,\\[1mm]
    x_3^{(k)}&=-\frac37\,x_1^{(k-1)}-\frac37\,x_2^{(k-1)}\,,
  \end{aligned}\right.
  \qquad \underline{k=1,2,\dots}
\]
\[ x_1^{(0)}=0,\quad x_2^{(0)}=0,\quad x_3^{(0)}=0 \]
$k=1$ olduğunda
\[
  \left\{\begin{aligned}
    x_1^{(1)}&=\frac13\,x_2^{(0)}-\frac13\,x_3^{(0)}+\frac13=\frac13\\[1mm]
    x_2^{(1)}&=-\frac12\,x_1^{(0)}-\frac13\,x_3^{(0)}=0\\[1mm]
    x_3^{(1)}&=-\frac37\,x_1^{(0)}-\frac37\,x_2^{(0)}=0
  \end{aligned}\right.
\]
\[ x_1^{(1)}=\frac13,\quad x_2^{(1)}=0,\quad x_3^{(1)}=0 \]
$k=2$ olduğunda
\[
  \left\{\begin{aligned}
    x_1^{(2)}&=\frac13\,x_2^{(1)}-\frac13\,x_3^{(1)}+\frac13=\frac13\\[1mm]
    x_2^{(2)}&=-\frac12\,x_1^{(1)}-\frac13\,x_3^{(1)}=-\frac16\\[1mm]
    x_3^{(2)}&=-\frac37\,x_1^{(1)}-\frac37\,x_2^{(1)}=-\frac37\cdot\frac13=-\frac17
  \end{aligned}\right.
\]
Böylece,
\[ x_1^{(2)}=\frac13,\quad x_2^{(2)}=-\frac16,\quad x_3^{(2)}=-\frac17 \]

\end{cozum}

\sayfa{11}
Şimdi Jacobi iterasyon formülünü matris-vektör şeklinde verelim. Bunun için $A$ matrisini
{\small\setlength{\arraycolsep}{2.5pt}
\[
  \underbrace{\begin{bmatrix}a_{11}&a_{12}&\cdots&a_{1n}\\a_{21}&a_{22}&\cdots&a_{2n}\\\vdots&&&\vdots\\a_{n1}&a_{n2}&\cdots&a_{nn}\end{bmatrix}}_{A}
  =\underbrace{\begin{bmatrix}0&0&\cdots&0\\a_{21}&0&\cdots&0\\\vdots&&&\vdots\\a_{n1}&a_{n2}&\cdots&0\end{bmatrix}}_{\substack{L\\ \text{(alt üçgen)}}}
  +\underbrace{\begin{bmatrix}a_{11}&0&\cdots&0\\0&a_{22}&\cdots&0\\\vdots&&&\vdots\\0&0&\cdots&a_{nn}\end{bmatrix}}_{\substack{D\\ \text{(köşegen)}}}
  +\underbrace{\begin{bmatrix}0&a_{12}&\cdots&a_{1n}\\0&0&\cdots&a_{2n}\\\vdots&&&\vdots\\0&0&\cdots&0\end{bmatrix}}_{\substack{U\\ \text{(üst üçgen)}}}
\]}
\[ A=L+D+U \]
Şeklinde yazalım. Bu durumda
\[ Ax=b \ \Leftrightarrow\ (L+D+U)\,x=b \ \Rightarrow\ Dx=-(L+U)\,x+b \]
olur. $a_{ii}\neq 0,\ i=1,\dots,n$ olduğuna göre $\det D\neq 0$ olur ve demek ki, tersi vardır: Her tarafı $D^{-1}$'le çarparsak
\[ x=-D^{-1}(L+U)\,x+D^{-1}b \ \Leftrightarrow\ \underline{x=T_J\,x+c_J} \]
elde edilir.
Başlangıç iterasyon olarak $x^{0}=0$ alırsak, bu denklem üzerinden
\[ T_J=-D^{-1}(L+U),\qquad c_J=D^{-1}b \]
işaretlersek
\formulkutu{x^{(k)}=T_J\,x^{(k-1)}+c_J,\quad k=1,2,\dots}
Şeklinde iterasyon belirleyebiliriz.

$T_J$ matrisine Jacobi matrisi denir.

\sayfa{12}
\underline{Örneğin} bir önceki örnekteki
\[
  \begin{cases}
    \kutula{3x_1}-x_2+x_3=1\\
    3x_1+\kutula{6x_2}+2x_3=0\\
    3x_1+3x_2+\kutula{7x_3}=0
  \end{cases}
  \Leftrightarrow\
  \underbrace{\begin{bmatrix}3&-1&1\\3&6&2\\3&3&7\end{bmatrix}}_{A}
  \underbrace{\begin{bmatrix}x_1\\x_2\\x_3\end{bmatrix}}_{x}
  =\underbrace{\begin{bmatrix}1\\0\\0\end{bmatrix}}_{b}
\]
Sistemi için
\[
  D=\begin{bmatrix}3&0&0\\0&6&0\\0&0&7\end{bmatrix}
  \quad\text{ve}\quad
  L+U=\begin{bmatrix}0&-1&1\\3&0&2\\3&3&0\end{bmatrix}
\]
\[
  \text{olduğundan ve}\quad
  D^{-1}=\begin{bmatrix}\frac13&0&0\\[1mm]0&\frac16&0\\[1mm]0&0&\frac17\end{bmatrix}
  \quad\text{olduğu için}
\]
\[ Ax=b \ \Leftrightarrow\ (D+L+U)\,x=b \ \Leftrightarrow\ Dx=-(L+U)\,x+b \]
\[ \Leftrightarrow\ x=\underbrace{-D^{-1}(L+U)}_{T_J}\,x+\underbrace{D^{-1}b}_{c_J} \]
\[
  T_J=-\begin{bmatrix}\frac13&0&0\\[1mm]0&\frac16&0\\[1mm]0&0&\frac17\end{bmatrix}\cdot\begin{bmatrix}0&-1&1\\3&0&2\\3&3&0\end{bmatrix}
  =-\begin{bmatrix}0&-\frac13&\frac13\\[1mm]\frac12&0&\frac13\\[1mm]\frac37&\frac37&0\end{bmatrix}
\]
\[
  \Rightarrow\quad \text{Jacobi matrisi}\quad
  T_J=\begin{bmatrix}0&\frac13&\frac13\\[1mm]-\frac12&0&-\frac13\\[1mm]-\frac37&-\frac37&0\end{bmatrix}\footnotemark
  \quad\text{olur.}
\]
\footnotetext{Claude notu: bir önceki satırdaki çarpıma göre birinci satır muhtemelen $\bigl[0,\ \tfrac13,\ -\tfrac13\bigr]$ olmalı.}

\sayfa{13}
\section{Sorular}
\begin{enumerate}[label=\arabic*)]
  \item Aşağıdaki matrislerin özdeğer ve özfonksiyonlarını bulunuz:
  \[
    \text{a)}\ A=\begin{bmatrix}1&2\\-1&2\end{bmatrix},\qquad
    \text{b)}\ A=\begin{bmatrix}-1&2&0\\0&3&4\\0&0&7\end{bmatrix}
  \]
  Ayrıca bu matrislerin spektral yarıçaplarını hesaplayınız.
  \item
  \[
    \begin{cases}
      4x_1+x_2-x_3+x_4=-2\\
      x_1+4x_2-x_3-x_4=-1\\
      -x_1-x_2+5x_3+x_4=0\\
      x_1-x_2+x_3+3x_4=1
    \end{cases}
  \]
  Sistemini $x^{(0)}=0$ olmakla Jacobi iterasyonu kullanarak ilk iki iterasyonunu bulunuz. Jacobi matrisini bulunuz.
\end{enumerate}

\end{document}
